% ===========================================================================
\chapter{Reading Datalog: the parser}\label{ch:parser}

Tokens say what the words are, not what they mean together. The rules for
how tokens combine are Datalog's \emph{grammar}, and they are short enough
to write down in full, as the first lines of the parser:

\begin{newfile}{parser.nex}
-- parser.nex: reading a Datalog program.
--
--   program  ::= { clause }
--   clause   ::= atom [ ":-" literal { "," literal } ] "."
--              | "?-" atom "."
--   literal  ::= atom | "not" atom | term op term
--   op       ::= "=" | "!=" | "<" | "<=" | ">" | ">="
--   atom     ::= name [ "(" term { "," term } ")" ]
--   term     ::= name | variable | integer | string

intern Lexer
intern data/Result

class Parser
create
  make(text: String) do
    let lexer := create Lexer.make(text)
    tokens := []
    lines := []
    from let done := false until done do
      let t := lexer.next_token
      tokens.add(t)
      lines.add(lexer.line)
      match t of
        End_Of_Input then done := true
        _            then done := false
      end
    end
    index := 0
    anonymous := 0
  end
feature
  tokens: Array[Token]
  -- The line each token is on.
  lines: Array[Integer]
  index: Integer
  -- How many `_` variables have been read, so that each gets its own name.
  anonymous: Integer
\end{newfile}

Read \code{::=} as ``is'', \code{|} as ``or'', braces as ``any number of
times'' and square brackets as ``optionally''. A program is any number of
clauses; a clause is a fact, a rule or a question; and so on down to a
term.

The parser reads the whole text into tokens first, recording each token's
line beside it, and then walks along the list. The grammar translates into
code almost word for word: one routine for each rule of the grammar, each
calling the routines for the rules it mentions. This style is called
\emph{recursive descent}.

\section{Moving along the tokens}

A few small routines do the walking. \code{advance} has the precondition
\code{not_at_end}: no correct parser ever steps past \code{End_Of_Input},
and if one did, the mistake would be reported at the step that went too
far. \code{located} prefixes a message with the line it is about, and
\code{expect} uses it when the punctuation it needs is not there.

\begin{continued}{parser.nex}
  current: Token do
    result := tokens.get(index)
  end

  at_end: Boolean do
    match current of
      End_Of_Input then result := true
      _            then result := false
    end
  end

  at_punct(p: String): Boolean do
    match current of
      Punct(text) then result := text = p
      _           then result := false
    end
  end

  advance
  require
    not_at_end: not at_end
  do
    index := index + 1
  ensure
    moved_on: index = old index + 1
  end

  -- `message`, saying where in the program it applies.
  located(message: String): String do
    result := "line " + lines.get(index).to_string + ": " + message
  end

  expect(p: String) do
    if at_punct(p) then
      advance
    else
      raise located("expected '" + p + "' but found " + token_text(current))
    end
  end
\end{continued}

\section{Clauses}

A program is clauses until the end of the input. A clause that starts with
\code{?-} is a question; otherwise it starts with an atom, which is the
head of a rule if \code{:-} follows and a fact if a full stop does.

\begin{continued}{parser.nex}
  -- The whole program.
  program: Datalog_Program do
    result := create Datalog_Program.make
    from until at_end do
      clause(result)
    end
  ensure
    all_read: at_end
  end

  clause(p: Datalog_Program) do
    if at_punct("?-") then
      advance
      p.add_query(atom)
      expect(".")
    else
      let head := atom
      if at_punct(":-") then
        advance
        let body: Array[Literal] := [literal]
        from until not at_punct(",") do
          advance
          body.add(literal)
        end
        expect(".")
        p.add_rule(create Rule.make(head, body))
      else
        expect(".")
        if not head.is_ground then
          raise located("a fact cannot have variables: " + head.text)
        end
        p.add_fact(head)
      end
    end
  end
\end{continued}

A fact with a variable in it is not a fact, so the parser refuses it with
a message, rather than leaving it to the precondition of \code{add_fact}.
That is the division of labour this tutorial follows throughout: mistakes
in the \emph{input} are reported as messages, and preconditions are kept
for mistakes in the \emph{program}.

\section{Literals}

A literal can begin in three ways. \code{not} followed by a name is a
negated atom. A name \emph{not} followed by a comparison operator is an
atom. Anything else is a test such as \code{X != Y}, whose left side may be
a constant as well as a variable. Telling these apart needs a peek at the
token after the current one, which \code{ahead} provides.

\begin{continued}{parser.nex}
  literal: Literal do
    if is_name("not", 0) and is_any_name(1) then
      advance
      result := create Negated.make(atom)
    elseif is_any_name(0) and not is_operator(1) then
      result := create Positive.make(atom)
    else
      let left := term
      if not is_operator(0) then
        raise located("expected a comparison but found " + token_text(current))
      end
      let op := punct_text(current)
      advance
      result := create Compare.make(op, left, term)
    end
  end

  -- The token `k` places after the current one, or the last token.
  ahead(k: Integer): Token do
    result := tokens.get((index + k).min(tokens.length - 1))
  end

  is_name(name: String, k: Integer): Boolean do
    match ahead(k) of
      Name_Token(text) then result := text = name
      _                then result := false
    end
  end

  is_any_name(k: Integer): Boolean do
    match ahead(k) of
      Name_Token then result := true
      _          then result := false
    end
  end

  is_operator(k: Integer): Boolean do
    result := ["=", "!=", "<", "<=", ">", ">="].contains(punct_text(ahead(k)))
  end
\end{continued}

\section{Atoms and terms}

An atom is a name, optionally followed by terms in parentheses. A term is a
single token, and its kind decides what it becomes. One variable is special:
\code{_} means ``something, I don't care what'', and each \code{_} must be
a different variable, or \code{flight(A, _, _)} would only match flights
whose second and third arguments were equal. So the parser gives each one
its own name, \code{_1}, \code{_2} and so on.

\begin{continued}{parser.nex}
  atom: Atom do
    match current of
      Name_Token(text) then
        do
          advance
          let args: Array[Term] := []
          if at_punct("(") then
            advance
            args.add(term)
            from until not at_punct(",") do
              advance
              args.add(term)
            end
            expect(")")
          end
          result := create Atom.make(text, args)
        end
      _ then raise located("expected a predicate but found " + token_text(current))
    end
  end

  term: Term do
    match current of
      Name_Token(text) then
        do
          advance
          result := create Symbol.make(text)
        end
      Var_Token(text) if text = "_" then
        do
          advance
          anonymous := anonymous + 1
          result := create Var.make("_" + anonymous.to_string)
        end
      Var_Token(text) then
        do
          advance
          result := create Var.make(text)
        end
      Int_Token(value) then
        do
          advance
          result := create Int.make(value)
        end
      Text_Token(value) then
        do
          advance
          result := create Text.make(value)
        end
      _ then raise located("expected a term but found " + token_text(current))
    end
  end
invariant
  same_length: lines.length = tokens.length
  inside_tokens: 0 <= index and index < tokens.length
end
\end{continued}

\section{Errors at the boundary}

The parser raises an exception at the first mistake. That is convenient
inside it: a mistake found five calls deep in \code{term} goes straight
out, without every routine in between checking and passing it along. But a
caller should not have to know that. Mistakes in a program are an ordinary
outcome, so the one function the rest of the engine uses turns them into a
\code{Result} from Nex's \code{data/Result} library: an \code{Ok} holding
the program, or an \code{Err} holding the message.

\begin{continued}{parser.nex}
-- The program `text` describes, or what is wrong with it.
function parse_program(text: String): Result[Datalog_Program, String] do
  let parser := create Parser.make(text)
  result := create Ok[Datalog_Program, String].make(parser.program)
rescue
  result := create Err[Datalog_Program, String].make(exception.to_string)
end
\end{continued}

\section{Trying it out}

\begin{tryit}{try_parser.nex}
intern Parser
intern io/Path

function show(text: String) do
  match parse_program(text) of
    Ok as ok then
      do
        across ok.value.facts as f do
          console.print_line("fact      " + f.text + ".")
        end
        across ok.value.rules as r do
          console.print_line("rule      " + r.text)
        end
        across ok.value.queries as q do
          console.print_line("question  " + q.text)
        end
      end
    Err as e then console.print_line("error: " + e.error)
  end
end

show(create Path.make(program.working_directory + "/family.dl").read_text)
show("airport(A) :- flight(A, _, _).")
show("parent(alice bob).")
show("parent(X, bob).")
show("p(a).\nq(b) :- p(a)\nr(c).")
show("p(a) & q(b).")
\end{tryit}
\tryoutput{try_parser}

The parser reads every clause of \code{family.dl} and prints it back. The
\code{_} variables have been numbered. The mistakes are reported with their
lines: the full stop missing at the end of line 2, after
\code{q(b) :- p(a)}, is only noticed on line 3, when \code{r} turns up
where the full stop should be.
